\chapter{Hiện thực}

\section{Mục tiêu của chương}

Chương này trình bày cách hệ thống được hiện thực dựa trên các quyết định thiết kế ở Chương 2 và Chương 3. Nội dung tập trung vào phương pháp hiện thực, công nghệ, tổ chức mã nguồn, backend, frontend, database, API, business rules, quản lý mã nguồn và kết quả hiện thực theo từng phase.

\section{Phương pháp hiện thực}

Nhóm hiện thực hệ thống theo từng module nghiệp vụ. Với mỗi module, nhóm triển khai theo thứ tự: thiết kế dữ liệu, backend API, frontend UI, tích hợp FE--BE, kiểm thử chức năng và ghi nhận bằng chứng. Cách làm này phù hợp với quy trình phát triển theo use case vì mỗi chức năng có thể truy vết từ yêu cầu đến API, database, giao diện và test.

\section{Cơ sở kiểm chứng codebase}

Do báo cáo cần phản ánh đúng hệ thống đã hiện thực, nhóm sử dụng codebase làm nguồn kiểm chứng chính khi mô tả kết quả hiện thực. Các số liệu như số lượng module, controller, service, endpoint, route frontend, migration và test chỉ nên được giữ trong báo cáo nếu đã được đối chiếu với mã nguồn tại thời điểm viết.

\begin{longtable}{|p{4cm}|p{6cm}|p{4.4cm}|}
\caption{Nguồn kiểm chứng nội dung hiện thực trong codebase}\\
\hline
\textbf{Hạng mục} & \textbf{Nguồn kiểm chứng} & \textbf{Mục đích kiểm tra} \\
\hline
Prisma models/enums/migrations & \path{backend/prisma/schema.prisma}, \path{backend/prisma/migrations} & Kiểm tra số lượng bảng, enum, quan hệ và lịch sử thay đổi schema. \\
\hline
Backend modules & \code{backend/src/modules} hoặc các thư mục module tương ứng & Kiểm tra module nghiệp vụ đã hiện thực. \\
\hline
Controllers/endpoints & \path{backend/src/**/*.controller.ts}, Swagger UI & Kiểm tra số lượng API, method, path và role/guard. \\
\hline
Services/business rules & \code{backend/src/**/*.service.ts} & Kiểm tra các rule thật sự được enforce ở service layer. \\
\hline
Frontend routes/pages & \code{frontend/src/app}, \code{frontend/src/features}, \code{frontend/src/config} & Kiểm tra màn hình, route, sidebar và phân quyền giao diện. \\
\hline
Test files & \code{backend/test}, \code{backend/src/**/*.spec.ts}, \code{frontend/src/**/*.test.tsx} nếu có & Kiểm tra phạm vi test tự động hiện có. \\
\hline
Package versions & \code{backend/package.json}, \code{frontend/package.json} & Kiểm tra phiên bản NestJS, React, Vite, Tailwind và các thư viện liên quan. \\
\hline
\end{longtable}

\noindent
Vì codebase có thể tiếp tục thay đổi trong quá trình hoàn thiện đồ án, các số liệu định lượng trong chương này như số endpoint, số route hoặc số test cần được xem là số liệu tại thời điểm audit. Nếu codebase thay đổi, nhóm cần cập nhật lại bảng tương ứng để tránh lệch giữa báo cáo và hiện thực.

\section{Quy trình hiện thực theo use case}

Nhóm hiện thực hệ thống theo hướng bám use case. Thay vì triển khai giao diện hoặc API rời rạc, mỗi chức năng được đi theo một chuỗi công việc từ yêu cầu đến kiểm thử. Cách làm này giúp mỗi phần code có thể truy vết ngược về yêu cầu và business rules.

\begin{longtable}{|p{1.2cm}|p{4cm}|p{9cm}|}
\caption{Quy trình hiện thực một chức năng nghiệp vụ}\\
\hline
\textbf{Bước} & \textbf{Hoạt động} & \textbf{Kết quả mong đợi} \\
\hline
1 & Đối chiếu yêu cầu và use case & Xác định actor, input, output, precondition, postcondition và business rules liên quan. \\
\hline
2 & Kiểm tra thiết kế dữ liệu & Xác định model, relation, enum và constraint cần dùng hoặc cần bổ sung. \\
\hline
3 & Cập nhật Prisma schema/migration nếu cần & Database schema phản ánh đúng nghiệp vụ cần hiện thực. \\
\hline
4 & Xây dựng DTO & Request body được validate rõ ràng ở biên API. \\
\hline
5 & Hiện thực service layer & Business rules, state transition và transaction được xử lý tại backend. \\
\hline
6 & Hiện thực controller/API & Endpoint REST được expose theo prefix \code{/api/v1} và có Swagger metadata nếu có. \\
\hline
7 & Gắn guard/role & API nghiệp vụ được bảo vệ bằng JWT và role phù hợp. \\
\hline
8 & Xây dựng frontend API client/hook & Frontend gọi API thông qua lớp client/hook thống nhất. \\
\hline
9 & Xây dựng giao diện & Người dùng thao tác được theo role, có loading/error/success state. \\
\hline
10 & Kiểm thử chức năng & Chạy E2E/API/manual test, kiểm tra cả luồng đúng và luồng lỗi. \\
\hline
11 & Ghi nhận bằng chứng & Lưu screenshot UI, Swagger response, test log hoặc đoạn code minh họa. \\
\hline
\end{longtable}

\section{Cấu trúc mã nguồn tổng quan}

\begin{lstlisting}[language=bash,caption={Cấu trúc codebase rút gọn}]
4N_Clinic_Management_System/
|-- backend/
|   |-- prisma/              # schema.prisma, migrations, seed.ts
|   |-- src/
|   |   |-- common/          # guards, decorators, filters
|   |   |-- prisma/          # PrismaService
|   |   `-- modules/         # 21 feature folders
|   `-- test/                # e2e tests
|-- frontend/
|   `-- src/
|       |-- app/             # router, providers
|       |-- components/      # common + ui components
|       |-- config/          # navigation, permissions
|       |-- features/        # feature-based pages
|       `-- lib/             # api-client, helpers
`-- docs/                    # design docs and evidence
\end{lstlisting}

\section{Hiện thực backend}

Backend sử dụng NestJS 11 và TypeScript. Codebase hiện có 21 feature folders, 20 controllers có API route và 24 service files (21 service nghiệp vụ + \code{PrismaService}, \code{AppService}, \code{HealthService}). Các module được chia theo domain nghiệp vụ thay vì chia theo kiểu kỹ thuật thuần túy. Ví dụ, module \code{visits} chịu trách nhiệm tiếp nhận lượt khám, module \code{examinations} chịu trách nhiệm phiếu khám và kê đơn, module \code{billing} chịu trách nhiệm hóa đơn và thanh toán.

\begin{longtable}{|p{3.2cm}|p{2cm}|p{2.5cm}|p{6cm}|}
\caption{Module backend theo phase}\\
\hline
\textbf{Module} & \textbf{Phase} & \textbf{Endpoint} & \textbf{Vai trò} \\
\hline
auth & P1 & 4 & Login, refresh, logout, me \\
\hline
users & P1 & 6 & Quản lý tài khoản \\
\hline
rbac & P1 & 3 & Vai trò và quyền \\
\hline
patients & P1 & 4 & Hồ sơ bệnh nhân, lịch sử khám \\
\hline
visits & P1 & 3 & Tiếp nhận lượt khám, mở examination \\
\hline
examinations & P1 & 6 & Phiếu khám, chẩn đoán, kê đơn, hoàn tất \\
\hline
prescriptions & P1 & 0 & Service-only module dùng bởi examinations \\
\hline
billing & P1 & 5 & Hóa đơn, chi tiết hóa đơn, thanh toán \\
\hline
diseases, drugs, regulations & P1 & 9 & Danh mục và quy định \\
\hline
reports & P1+P2 & 2 & Báo cáo tháng và revenue breakdown \\
\hline
appointments, queue, vitals & P2 & 12 & Lịch hẹn, hàng đợi, sinh hiệu \\
\hline
services, lab & P2 & 14 & Dịch vụ và xét nghiệm \\
\hline
inventory, pharmacy, organization, audit & P2 & 25 & Kho thuốc, cấp phát, tổ chức và audit \\
\hline
\end{longtable}

\section{Hiện thực frontend}

Frontend sử dụng React 19, Vite 7, TypeScript, Tailwind CSS 4, TanStack Query, Zustand, React Hook Form và Zod. Codebase hiện có 36 named routes và 33 page files. Các route được bảo vệ bằng \code{ProtectedRoute} và \code{RequireRole}. Sidebar được cấu hình theo role trong navigation config.

\begin{lstlisting}[language=bash,caption={Cấu trúc frontend rút gọn}]
frontend/src/
|-- app/
|   |-- router.tsx
|   `-- providers.tsx
|-- config/
|   |-- navigation.ts
|   `-- permissions.ts
|-- features/
|   |-- auth/
|   |-- patients/
|   |-- visits/
|   |-- examinations/
|   |-- invoices/
|   |-- reports/
|   |-- appointments/
|   |-- lab/
|   |-- inventory/
|   |-- pharmacy/
|   `-- audit/
`-- lib/
    |-- api-client.ts
    `-- theme.ts
\end{lstlisting}

\section{Hiện thực cơ sở dữ liệu}

Database được định nghĩa bằng Prisma schema và triển khai trên PostgreSQL. Prisma migration quản lý thay đổi schema theo từng giai đoạn: baseline, identity-access và phase2a-foundation. Seed data tạo dữ liệu demo cho tài khoản, danh mục, bệnh nhân, lượt khám, hóa đơn, dịch vụ, lô thuốc và các dữ liệu cần thiết cho demo.

\section{Hiện thực API}

API được thiết kế theo REST với prefix \code{/api/v1}. Hệ thống sử dụng Swagger UI tại \code{http://localhost:3000/api/docs} để kiểm tra và tài liệu hóa endpoint. Tất cả route nghiệp vụ đều yêu cầu JWT, ngoại trừ endpoint login.

\begin{table}[H]
\centering
\caption{Kết quả hiện thực API}
\begin{tabularx}{\textwidth}{|p{4cm}|X|}
\hline
\textbf{Hạng mục} & \textbf{Kết quả} \\
\hline
API prefix & \code{/api/v1} \\
\hline
Swagger UI & \code{http://localhost:3000/api/docs} \\
\hline
Tổng endpoints & 92 endpoints \\
\hline
Phase 1 endpoints & 41 endpoints \\
\hline
Phase 2 endpoints & 51 endpoints \\
\hline
Auth & JWT Access Token + Refresh Token \\
\hline
API protection & JwtAuthGuard, RolesGuard, @Roles decorator \\
\hline
\end{tabularx}
\end{table}

\section{Hiện thực business rules tiêu biểu}

Một điểm quan trọng của hệ thống là các quy tắc nghiệp vụ được enforce tại service layer. Một số nghiệp vụ có sử dụng \code{prisma.$transaction()$} để đảm bảo tính nguyên tử, ví dụ: tạo lượt khám và số thứ tự, mở phiên khám, thanh toán hóa đơn, kích hoạt quy định và cấp phát thuốc.

\begin{table}[H]
\centering
\caption{Transaction usage trong backend}
\begin{tabularx}{\textwidth}{|p{4cm}|p{3cm}|X|}
\hline
\textbf{Operation} & \textbf{Module} & \textbf{Mục đích} \\
\hline
Visit creation + queue number & visits & Tạo lượt khám và số thứ tự an toàn, tránh race condition \\
\hline
Open examination & visits & Cập nhật visit status và tạo examination atomically \\
\hline
Payment + invoice status & billing & Cập nhật paidAmount và InvoiceStatus cùng lúc \\
\hline
Regulation activation & regulations & Deactivate phiên bản cũ và activate phiên bản mới \\
\hline
Dispense + stock deduction & pharmacy & Tạo phiếu cấp phát, stock movement và giảm tồn kho \\
\hline
\end{tabularx}
\end{table}

\subsection{Ví dụ hiện thực business rule bằng transaction}

Đoạn mã dưới đây minh họa cách service layer xử lý BR-05 và BR-07 trong luồng tạo lượt khám. Code được rút gọn để tập trung vào ý tưởng chính: kiểm tra trùng lượt khám trong ngày, lấy số thứ tự cuối cùng và tạo visit trong cùng một transaction.

\begin{lstlisting}[language=JavaScript,caption={BR-05/BR-07: Tạo lượt khám và sinh số thứ tự trong Prisma transaction}]
async createVisit(dto, currentUser) {
  return this.prisma.$transaction(async (tx) => {
    const patient = await tx.patient.findUnique({
      where: { id: dto.patientId },
    });
    if (!patient) {
      throw new NotFoundException('Patient not found');
    }

    const duplicated = await tx.visit.findFirst({
      where: {
        patientId: dto.patientId,
        visitDate: dto.visitDate,
      },
    });
    if (duplicated) {
      throw new ConflictException('Patient already has a visit today');
    }

    const lastVisit = await tx.visit.findFirst({
      where: { visitDate: dto.visitDate },
      orderBy: { queueNumber: 'desc' },
    });
    const queueNumber = (lastVisit?.queueNumber ?? 0) + 1;

    return tx.visit.create({
      data: {
        patientId: dto.patientId,
        visitDate: dto.visitDate,
        queueNumber,
        createdById: currentUser.id,
        status: 'WAITING',
      },
    });
  });
}
\end{lstlisting}

Mẫu hiện thực này thể hiện nguyên tắc quan trọng của hệ thống: frontend có thể hỗ trợ validation ban đầu, nhưng các business rules bắt buộc vẫn phải được kiểm tra tại backend service và database transaction.

\section{Quản lý mã nguồn và quy trình nhóm}

Git history cho thấy nhóm sử dụng các branch theo convention \code{feature/UCxx-yy}, \code{docs/}, \code{db/}, \code{chore/}. Commit timeline phản ánh quá trình từ khởi tạo repo, setup backend/frontend, triển khai Phase 1, chốt baseline v1, mở rộng Phase 2 và bổ sung seed data.

\begin{longtable}{|p{3.2cm}|p{3cm}|p{7.5cm}|}
\caption{Timeline phát triển từ git history}\\
\hline
\textbf{Giai đoạn} & \textbf{Thời gian} & \textbf{Nội dung} \\
\hline
Khởi tạo & 2026-03-12 đến 2026-03-14 & Initial commit, cấu trúc ban đầu \\
\hline
Setup project & 2026-04-19 & Cấu trúc backend, docs \\
\hline
Phase 1 backend core & 2026-04-27 & Prisma schema, Phase 1 modules, seed \\
\hline
Phase 1 UC07--UC11 & 2026-05-17 & Visit và examination flow \\
\hline
Phase 1 UC01--UC03 + frontend & 2026-05-19 & Auth, users, RBAC và merge \\
\hline
Phase 1 baseline & 2026-05-19 & README v1, baseline đầu tiên \\
\hline
Phase 2 full implementation & 2026-06-05 & Backend và frontend Phase 2 \\
\hline
Bug fixes + seed & 2026-06-05 đến 2026-06-06 & Fix examination, comprehensive seed data \\
\hline
\end{longtable}

\textbf{Lưu ý về git history:} Git log phản ánh hoạt động commit, không phản ánh toàn bộ công sức đóng góp. Nhóm có thực hiện pair programming, review thiết kế trực tiếp, kiểm thử thủ công và chuẩn bị tài liệu ngoài commit. Nhóm xác nhận 4 thành viên đóng góp tương đương; bằng chứng đầy đủ gồm Git history, biên bản phân công, checklist kiểm thử và screenshot demo.

\section{Demo các chức năng trọng tâm}

\begin{longtable}{|p{3.2cm}|p{5.5cm}|p{5.5cm}|}
\caption{Demo hiện thực chức năng}\\
\hline
\textbf{Demo} & \textbf{Nội dung trình bày} & \textbf{Minh chứng } \\
\hline
Demo Auth/RBAC & Đăng nhập, lấy thông tin user, phân quyền API & Swagger response, screenshot sidebar \\
\hline
Demo Patients & Tạo bệnh nhân mới, tìm kiếm bệnh nhân & Screenshot UI, API response \\
\hline
Demo Visits & Tạo lượt khám, cấp số thứ tự, chặn trùng ngày & Screenshot, đoạn service rule \\
\hline
Demo Examinations & Ghi phiếu khám, chẩn đoán, kê đơn & Screenshot phiếu khám \\
\hline
Demo Billing & Lập hóa đơn, thanh toán nhiều lần & Screenshot hóa đơn, payment dialog \\
\hline
Demo Phase 2 & Lịch hẹn, xét nghiệm, kho thuốc, FEFO & Lab/pharmacy/inventory screenshots \\
\hline
\end{longtable}

\ReportFigure{figures/screenshots/pharmacy-worklist.png}{Giao diện Pharmacy Worklist -- minh chứng hiện thực Phase 2}{fig:implementation-demo}

\section{Kết luận chương}

Hệ thống đã được chuyển hoá từ thiết kế thành sản phẩm chạy được qua hai phase. Phase 1 hoàn thành baseline nghiệp vụ lõi gồm auth/RBAC, quản lý bệnh nhân, lượt khám, phiếu khám, kê đơn, hóa đơn, thanh toán, danh mục và báo cáo tháng. Phase 2 mở rộng với lịch hẹn, hàng đợi, sinh hiệu, dịch vụ, xét nghiệm, kho thuốc, cấp phát thuốc và audit log. API và business rules đều được enforce tại backend service layer; giao diện theo vai trò được bảo vệ bởi ProtectedRoute và RequireRole.
